%% ECON 282E -- Session 5, Deck B6: Worked Example -- the RBC Model by Projection
%% The instructor's request: Block A solves the quarterly RBC by perturbation,
%% Block B solves the SAME model by projection, and the session closes with one
%% model and four methods in one table.
%% Numbers and coordinates from tools/figures/s05_numbers.json, keys
%% "worked_proj", "three_methods", "fig_proj_policy".

%% ECON 282E -- Foundations of Macroeconomics (UC Riverside, Fall 2026)
%% Shared Beamer preamble for every deck in slides/S01 ... slides/S10.
%%
%% Usage (a deck lives in slides/S0X/ and is compiled from that folder):
%%   %% ECON 282E -- Foundations of Macroeconomics (UC Riverside, Fall 2026)
%% Shared Beamer preamble for every deck in slides/S01 ... slides/S10.
%%
%% Usage (a deck lives in slides/S0X/ and is compiled from that folder):
%%   %% ECON 282E -- Foundations of Macroeconomics (UC Riverside, Fall 2026)
%% Shared Beamer preamble for every deck in slides/S01 ... slides/S10.
%%
%% Usage (a deck lives in slides/S0X/ and is compiled from that folder):
%%   \input{../preamble/course_preamble.tex}
%%   \decktitle{1}{A1}{Who I Am \& Why This Course}{September 24, 2026}
%%   \begin{document}
%%   \makedecktitle
%%   ...
%%
%% Adapted from Courses/AI-ECON-2026/source/shared_preamble.tex (Metropolis theme,
%% concept/caution/example/takeaway boxes, \sectiondivider). Changes for this course:
%%   * course title block (\decktitle / \makedecktitle) instead of \makebeamertitle
%%   * \zh{...} typesets Chinese (book titles) under pdflatex via CJKutf8
%%   * researchbox for the "Research in the wild" frames
%%   * section pages off; TikZ libraries and the course map loaded here
%% Build with pdflatex only (two passes); tools/build_slides.py does this for you.

\documentclass[10pt,english,aspectratio=169]{beamer}

\usepackage[T1]{fontenc}
\usepackage{textcomp}
\usepackage[utf8]{inputenc}
\usepackage{babel}
\usepackage{CJKutf8}

\usepackage{amsmath, amssymb, amsthm, graphicx}
\usepackage{booktabs, multirow, tabularx}
\usepackage{tikz}
\usetikzlibrary{positioning,arrows.meta,shapes,calc,fit,backgrounds}
\usepackage{pgfplots}
\pgfplotsset{compat=1.16}
\usepackage[most]{tcolorbox}
\tcbuselibrary{skins, breakable}
\usepackage{listings}
\usepackage{xcolor}

\lstset{
  basicstyle=\ttfamily\small,
  keywordstyle=\color{blue!80!black}\bfseries,
  commentstyle=\color{green!50!black}\itshape,
  stringstyle=\color{orange!80!black},
  backgroundcolor=\color{gray!8},
  frame=single,
  rulecolor=\color{gray!40},
  breaklines=true,
  showstringspaces=false,
  numbers=none,
}

\usetheme{metropolis}
\usefonttheme{professionalfonts}
\metroset{sectionpage=none}

\definecolor{LightGray}{RGB}{230,230,230}
\definecolor{RedTitle}{RGB}{0,0,200}
\definecolor{VeryLightGray}{RGB}{245,245,245}
\definecolor{DarkText}{RGB}{0,0,0}
\definecolor{UCRBlue}{RGB}{0,45,106}
\definecolor{UCRGold}{RGB}{241,171,0}

% Stage colours of the course arc (used by the course map and elsewhere)
\colorlet{StageTrunk}{blue!12}
\colorlet{StageBranch}{cyan!14}
\colorlet{StageMachine}{gray!18}
\colorlet{StageWall}{red!14}
\colorlet{StageTools}{orange!20}
\colorlet{StageData}{green!16}
\colorlet{StageAgents}{violet!14}

\hypersetup{bookmarks=false,breaklinks=true,pdfborder={0 0 0},colorlinks=true,
  linkcolor=DarkText,urlcolor=RedTitle}

\setbeamercolor{background canvas}{bg=white}
\setbeamercolor{title}{fg=RedTitle}
\setbeamercolor{frametitle}{bg=VeryLightGray,fg=DarkText}
\setbeamercolor{normal text}{fg=DarkText}
\setbeamercolor{alerted text}{fg=RedTitle}
\setbeamersize{text margin left=5mm, text margin right=5mm}
\addtobeamertemplate{frametitle}{}{\vspace{-0.5em}}

\setbeamertemplate{navigation symbols}{}
\setbeamertemplate{footline}{%
  \hspace{5mm}{\usebeamerfont{page number in head/foot}\color{black!45}%
  ECON 282E \,$\cdot$\, \insertshortdeck}\hfill%
  \usebeamercolor[fg]{page number in head/foot}%
  \usebeamerfont{page number in head/foot}%
  \insertframenumber\,/\,\inserttotalframenumber\kern1em\vskip2pt%
}

% ---------------------------------------------------------------------------
% Chinese text under pdflatex (book titles etc.)
\newcommand{\zh}[1]{\begin{CJK*}{UTF8}{gbsn}#1\end{CJK*}}

% ---------------------------------------------------------------------------
% Deck title block
%   \decktitle{session}{deck id}{title}{date}
\newcommand{\insertshortdeck}{}
\newcommand{\decksession}{}
\newcommand{\deckid}{}
\newcommand{\decktitletext}{}
\newcommand{\deckdate}{}
\newcommand{\decktitle}[4]{%
  \renewcommand{\decksession}{#1}%
  \renewcommand{\deckid}{#2}%
  \renewcommand{\decktitletext}{#3}%
  \renewcommand{\deckdate}{#4}%
  \renewcommand{\insertshortdeck}{Session #1 \,$\cdot$\, Deck #1.#2}%
  \title{#3}%
  \author{Zhigang Feng}%
  \date{#4}%
}
\newcommand{\makedecktitle}{%
  \begin{frame}[plain,noframenumbering]
    \begin{tikzpicture}[remember picture,overlay]
      \fill[UCRBlue] (current page.north west) rectangle ([yshift=-0.55cm]current page.north east);
      \fill[UCRGold] ([yshift=-0.55cm]current page.north west) rectangle ([yshift=-0.62cm]current page.north east);
      \node[anchor=west, text=white, font=\small\bfseries] at ([xshift=5mm,yshift=-0.275cm]current page.north west)
        {ECON 282E \,$\cdot$\, Foundations of Macroeconomics};
      \node[anchor=east, text=white, font=\small] at ([xshift=-5mm,yshift=-0.275cm]current page.north east)
        {UC Riverside \,$\cdot$\, Fall 2026};
    \end{tikzpicture}
    \vspace*{1.1cm}
    {\usebeamerfont{title}\color{black!55}\large Session \decksession\ \,$\cdot$\, Deck \decksession.\deckid\par}
    \vspace{0.25cm}
    {\usebeamerfont{title}\color{RedTitle}\LARGE\bfseries \decktitletext\par}
    \vspace{0.7cm}
    {\large Zhigang Feng\par}
    \vspace{0.15cm}
    {\color{black!65}\deckdate\par}
    \vfill
    {\footnotesize\itshape\color{black!55}Build it on paper $\;\to\;$ solve it on the machine $\;\to\;$ push it to the frontier with AI\par}
  \end{frame}}

% ---------------------------------------------------------------------------
% Section divider (plain frame, not counted)
\newcommand{\sectiondivider}[2]{%
\begin{frame}[plain,noframenumbering]
\begin{center}
\vfill
{\Large\textbf{#1}\par}
\vspace{0.35cm}
{\normalsize #2\par}
\vfill
\end{center}
\end{frame}
}

% ---------------------------------------------------------------------------
% Boxes
\newtcolorbox{conceptbox}[1][]{colback=blue!3, colframe=RedTitle!55, boxrule=0.35mm,
  arc=1mm, left=2mm, right=2mm, top=1mm, bottom=1mm, #1}
\newtcolorbox{cautionbox}[1][]{colback=red!3, colframe=red!50!black, boxrule=0.35mm,
  arc=1mm, left=2mm, right=2mm, top=1mm, bottom=1mm, #1}
\newtcolorbox{examplebox}[1][]{colback=green!3, colframe=green!45!black, boxrule=0.35mm,
  arc=1mm, left=2mm, right=2mm, top=1mm, bottom=1mm, #1}
\newtcolorbox{takeawaybox}[1][]{colback=VeryLightGray, colframe=RedTitle!55, boxrule=0.35mm,
  arc=1mm, left=2mm, right=2mm, top=1mm, bottom=1mm, #1}
\newtcolorbox{researchbox}[1][]{enhanced, colback=violet!4, colframe=violet!60!black,
  boxrule=0.35mm, arc=1mm, left=2mm, right=2mm, top=1mm, bottom=1mm,
  fonttitle=\bfseries\small, title={Research in the wild}, #1}

\newcommand{\compacttables}{\renewcommand{\arraystretch}{1.15}}
\newcommand{\src}[1]{{\tiny\color{black!55}#1}}

% ---------------------------------------------------------------------------
\input{../preamble/course_map.tex}

%%   \decktitle{1}{A1}{Who I Am \& Why This Course}{September 24, 2026}
%%   \begin{document}
%%   \makedecktitle
%%   ...
%%
%% Adapted from Courses/AI-ECON-2026/source/shared_preamble.tex (Metropolis theme,
%% concept/caution/example/takeaway boxes, \sectiondivider). Changes for this course:
%%   * course title block (\decktitle / \makedecktitle) instead of \makebeamertitle
%%   * \zh{...} typesets Chinese (book titles) under pdflatex via CJKutf8
%%   * researchbox for the "Research in the wild" frames
%%   * section pages off; TikZ libraries and the course map loaded here
%% Build with pdflatex only (two passes); tools/build_slides.py does this for you.

\documentclass[10pt,english,aspectratio=169]{beamer}

\usepackage[T1]{fontenc}
\usepackage{textcomp}
\usepackage[utf8]{inputenc}
\usepackage{babel}
\usepackage{CJKutf8}

\usepackage{amsmath, amssymb, amsthm, graphicx}
\usepackage{booktabs, multirow, tabularx}
\usepackage{tikz}
\usetikzlibrary{positioning,arrows.meta,shapes,calc,fit,backgrounds}
\usepackage{pgfplots}
\pgfplotsset{compat=1.16}
\usepackage[most]{tcolorbox}
\tcbuselibrary{skins, breakable}
\usepackage{listings}
\usepackage{xcolor}

\lstset{
  basicstyle=\ttfamily\small,
  keywordstyle=\color{blue!80!black}\bfseries,
  commentstyle=\color{green!50!black}\itshape,
  stringstyle=\color{orange!80!black},
  backgroundcolor=\color{gray!8},
  frame=single,
  rulecolor=\color{gray!40},
  breaklines=true,
  showstringspaces=false,
  numbers=none,
}

\usetheme{metropolis}
\usefonttheme{professionalfonts}
\metroset{sectionpage=none}

\definecolor{LightGray}{RGB}{230,230,230}
\definecolor{RedTitle}{RGB}{0,0,200}
\definecolor{VeryLightGray}{RGB}{245,245,245}
\definecolor{DarkText}{RGB}{0,0,0}
\definecolor{UCRBlue}{RGB}{0,45,106}
\definecolor{UCRGold}{RGB}{241,171,0}

% Stage colours of the course arc (used by the course map and elsewhere)
\colorlet{StageTrunk}{blue!12}
\colorlet{StageBranch}{cyan!14}
\colorlet{StageMachine}{gray!18}
\colorlet{StageWall}{red!14}
\colorlet{StageTools}{orange!20}
\colorlet{StageData}{green!16}
\colorlet{StageAgents}{violet!14}

\hypersetup{bookmarks=false,breaklinks=true,pdfborder={0 0 0},colorlinks=true,
  linkcolor=DarkText,urlcolor=RedTitle}

\setbeamercolor{background canvas}{bg=white}
\setbeamercolor{title}{fg=RedTitle}
\setbeamercolor{frametitle}{bg=VeryLightGray,fg=DarkText}
\setbeamercolor{normal text}{fg=DarkText}
\setbeamercolor{alerted text}{fg=RedTitle}
\setbeamersize{text margin left=5mm, text margin right=5mm}
\addtobeamertemplate{frametitle}{}{\vspace{-0.5em}}

\setbeamertemplate{navigation symbols}{}
\setbeamertemplate{footline}{%
  \hspace{5mm}{\usebeamerfont{page number in head/foot}\color{black!45}%
  ECON 282E \,$\cdot$\, \insertshortdeck}\hfill%
  \usebeamercolor[fg]{page number in head/foot}%
  \usebeamerfont{page number in head/foot}%
  \insertframenumber\,/\,\inserttotalframenumber\kern1em\vskip2pt%
}

% ---------------------------------------------------------------------------
% Chinese text under pdflatex (book titles etc.)
\newcommand{\zh}[1]{\begin{CJK*}{UTF8}{gbsn}#1\end{CJK*}}

% ---------------------------------------------------------------------------
% Deck title block
%   \decktitle{session}{deck id}{title}{date}
\newcommand{\insertshortdeck}{}
\newcommand{\decksession}{}
\newcommand{\deckid}{}
\newcommand{\decktitletext}{}
\newcommand{\deckdate}{}
\newcommand{\decktitle}[4]{%
  \renewcommand{\decksession}{#1}%
  \renewcommand{\deckid}{#2}%
  \renewcommand{\decktitletext}{#3}%
  \renewcommand{\deckdate}{#4}%
  \renewcommand{\insertshortdeck}{Session #1 \,$\cdot$\, Deck #1.#2}%
  \title{#3}%
  \author{Zhigang Feng}%
  \date{#4}%
}
\newcommand{\makedecktitle}{%
  \begin{frame}[plain,noframenumbering]
    \begin{tikzpicture}[remember picture,overlay]
      \fill[UCRBlue] (current page.north west) rectangle ([yshift=-0.55cm]current page.north east);
      \fill[UCRGold] ([yshift=-0.55cm]current page.north west) rectangle ([yshift=-0.62cm]current page.north east);
      \node[anchor=west, text=white, font=\small\bfseries] at ([xshift=5mm,yshift=-0.275cm]current page.north west)
        {ECON 282E \,$\cdot$\, Foundations of Macroeconomics};
      \node[anchor=east, text=white, font=\small] at ([xshift=-5mm,yshift=-0.275cm]current page.north east)
        {UC Riverside \,$\cdot$\, Fall 2026};
    \end{tikzpicture}
    \vspace*{1.1cm}
    {\usebeamerfont{title}\color{black!55}\large Session \decksession\ \,$\cdot$\, Deck \decksession.\deckid\par}
    \vspace{0.25cm}
    {\usebeamerfont{title}\color{RedTitle}\LARGE\bfseries \decktitletext\par}
    \vspace{0.7cm}
    {\large Zhigang Feng\par}
    \vspace{0.15cm}
    {\color{black!65}\deckdate\par}
    \vfill
    {\footnotesize\itshape\color{black!55}Build it on paper $\;\to\;$ solve it on the machine $\;\to\;$ push it to the frontier with AI\par}
  \end{frame}}

% ---------------------------------------------------------------------------
% Section divider (plain frame, not counted)
\newcommand{\sectiondivider}[2]{%
\begin{frame}[plain,noframenumbering]
\begin{center}
\vfill
{\Large\textbf{#1}\par}
\vspace{0.35cm}
{\normalsize #2\par}
\vfill
\end{center}
\end{frame}
}

% ---------------------------------------------------------------------------
% Boxes
\newtcolorbox{conceptbox}[1][]{colback=blue!3, colframe=RedTitle!55, boxrule=0.35mm,
  arc=1mm, left=2mm, right=2mm, top=1mm, bottom=1mm, #1}
\newtcolorbox{cautionbox}[1][]{colback=red!3, colframe=red!50!black, boxrule=0.35mm,
  arc=1mm, left=2mm, right=2mm, top=1mm, bottom=1mm, #1}
\newtcolorbox{examplebox}[1][]{colback=green!3, colframe=green!45!black, boxrule=0.35mm,
  arc=1mm, left=2mm, right=2mm, top=1mm, bottom=1mm, #1}
\newtcolorbox{takeawaybox}[1][]{colback=VeryLightGray, colframe=RedTitle!55, boxrule=0.35mm,
  arc=1mm, left=2mm, right=2mm, top=1mm, bottom=1mm, #1}
\newtcolorbox{researchbox}[1][]{enhanced, colback=violet!4, colframe=violet!60!black,
  boxrule=0.35mm, arc=1mm, left=2mm, right=2mm, top=1mm, bottom=1mm,
  fonttitle=\bfseries\small, title={Research in the wild}, #1}

\newcommand{\compacttables}{\renewcommand{\arraystretch}{1.15}}
\newcommand{\src}[1]{{\tiny\color{black!55}#1}}

% ---------------------------------------------------------------------------
%% Course map: the recurring "you are here" slide for ECON 282E.
%% Loaded by course_preamble.tex. Pattern follows AI-ECON-2026 lec01_map.tex, but the map
%% shows the whole ten-session course rather than one lecture.
%%
%%   \CourseMap{s}{label}   s = 1..10 highlights session s and prints "you are here: label"
%%                          (e.g. \CourseMap{1}{1.A2}); earlier sessions get a check mark.
%%   \CourseMap{0}{}        full overview, nothing highlighted.
%%
%% Note: TikZ option lists are comma-split before TeX conditionals run, so every \ifnum
%% below sits outside [...] option lists.

\newcommand{\CourseMap}[2]{%
\begin{center}
\begin{tikzpicture}[
    sess/.style={rectangle, rounded corners=2pt, draw=black!45, minimum width=1.34cm,
        minimum height=1.12cm, text width=1.26cm, align=center, inner sep=1pt},
    stage/.style={font=\tiny\bfseries, text=black!70},
]
  \foreach \i/\d/\t/\c in {%
      1/Sep 24/Intro \\ Growth/StageTrunk,
      2/Oct 1/HA $\cdot$ OLG \\ Policy/StageBranch,
      3/Oct 8/Num.\ DP \\ Python/StageMachine,
      4/Oct 15/Numerics \\ I/StageMachine,
      5/Oct 22/Numerics \\ II/StageMachine,
      6/Oct 29/The wall \\ \textbf{Midterm}/StageWall,
      7/Nov 5/ML \\ Deep nets/StageTools,
      8/Nov 12/RL \\ HA $+$ DL/StageTools,
      9/Nov 19/Text \\ LLMs/StageData,
      10/Dec 3/Agents \\ \textbf{Projects}/StageAgents} {
    \node[sess, fill=\c] (n\i) at ({1.46*(\i-1)}, 0)
      {{\scriptsize\bfseries S\i}\\[-1pt]{\tiny\color{black!60}\d}\\[1pt]{\tiny \t}};
    \ifnum\i<#1
      \node[font=\scriptsize, text=green!45!black] at ([xshift=-2.5mm,yshift=-2.2mm]n\i.north east) {$\checkmark$};
    \fi
    \ifnum\i=#1
      \draw[RedTitle, very thick, rounded corners=2pt]
        ([xshift=-0.6mm,yshift=-0.6mm]n\i.south west) rectangle ([xshift=0.6mm,yshift=0.6mm]n\i.north east);
      % keep the label inside the picture at both ends of the row
      \ifnum\i<3
        \node[font=\scriptsize\bfseries, text=RedTitle, anchor=south west] at ([yshift=1.2mm]n\i.north west)
          {$\blacktriangledown$\,you are here: #2};
      \else\ifnum\i>8
        \node[font=\scriptsize\bfseries, text=RedTitle, anchor=south east] at ([yshift=1.2mm]n\i.north east)
          {you are here: #2\,$\blacktriangledown$};
      \else
        \node[font=\scriptsize\bfseries, text=RedTitle, anchor=south] at ([yshift=1.2mm]n\i.north)
          {you are here: #2\,$\blacktriangledown$};
      \fi\fi
    \fi
  }
  % stage band under the sessions
  \foreach \a/\b/\lab/\c in {1/1/trunk/StageTrunk, 2/2/branches/StageBranch,
      3/5/the machine $\cdot$ classical methods/StageMachine, 6/6/the wall/StageWall,
      7/8/new tools/StageTools, 9/9/new data/StageData, 10/10/colleagues/StageAgents} {
    \fill[\c] ([yshift=-1.5mm]n\a.south west) rectangle ([yshift=-4.6mm]n\b.south east);
    \path ([yshift=-1.5mm]n\a.south west) -- ([yshift=-4.6mm]n\b.south east)
      node[midway, stage] {\lab};
  }
  \node[font=\footnotesize\itshape, text=black!70] at ([yshift=-9.5mm]$(n1.south)!0.5!(n10.south)$)
    {Build it on paper $\;\to\;$ solve it on the machine $\;\to\;$ push it to the frontier with AI};
\end{tikzpicture}
\end{center}
}


%%   \decktitle{1}{A1}{Who I Am \& Why This Course}{September 24, 2026}
%%   \begin{document}
%%   \makedecktitle
%%   ...
%%
%% Adapted from Courses/AI-ECON-2026/source/shared_preamble.tex (Metropolis theme,
%% concept/caution/example/takeaway boxes, \sectiondivider). Changes for this course:
%%   * course title block (\decktitle / \makedecktitle) instead of \makebeamertitle
%%   * \zh{...} typesets Chinese (book titles) under pdflatex via CJKutf8
%%   * researchbox for the "Research in the wild" frames
%%   * section pages off; TikZ libraries and the course map loaded here
%% Build with pdflatex only (two passes); tools/build_slides.py does this for you.

\documentclass[10pt,english,aspectratio=169]{beamer}

\usepackage[T1]{fontenc}
\usepackage{textcomp}
\usepackage[utf8]{inputenc}
\usepackage{babel}
\usepackage{CJKutf8}

\usepackage{amsmath, amssymb, amsthm, graphicx}
\usepackage{booktabs, multirow, tabularx}
\usepackage{tikz}
\usetikzlibrary{positioning,arrows.meta,shapes,calc,fit,backgrounds}
\usepackage{pgfplots}
\pgfplotsset{compat=1.16}
\usepackage[most]{tcolorbox}
\tcbuselibrary{skins, breakable}
\usepackage{listings}
\usepackage{xcolor}

\lstset{
  basicstyle=\ttfamily\small,
  keywordstyle=\color{blue!80!black}\bfseries,
  commentstyle=\color{green!50!black}\itshape,
  stringstyle=\color{orange!80!black},
  backgroundcolor=\color{gray!8},
  frame=single,
  rulecolor=\color{gray!40},
  breaklines=true,
  showstringspaces=false,
  numbers=none,
}

\usetheme{metropolis}
\usefonttheme{professionalfonts}
\metroset{sectionpage=none}

\definecolor{LightGray}{RGB}{230,230,230}
\definecolor{RedTitle}{RGB}{0,0,200}
\definecolor{VeryLightGray}{RGB}{245,245,245}
\definecolor{DarkText}{RGB}{0,0,0}
\definecolor{UCRBlue}{RGB}{0,45,106}
\definecolor{UCRGold}{RGB}{241,171,0}

% Stage colours of the course arc (used by the course map and elsewhere)
\colorlet{StageTrunk}{blue!12}
\colorlet{StageBranch}{cyan!14}
\colorlet{StageMachine}{gray!18}
\colorlet{StageWall}{red!14}
\colorlet{StageTools}{orange!20}
\colorlet{StageData}{green!16}
\colorlet{StageAgents}{violet!14}

\hypersetup{bookmarks=false,breaklinks=true,pdfborder={0 0 0},colorlinks=true,
  linkcolor=DarkText,urlcolor=RedTitle}

\setbeamercolor{background canvas}{bg=white}
\setbeamercolor{title}{fg=RedTitle}
\setbeamercolor{frametitle}{bg=VeryLightGray,fg=DarkText}
\setbeamercolor{normal text}{fg=DarkText}
\setbeamercolor{alerted text}{fg=RedTitle}
\setbeamersize{text margin left=5mm, text margin right=5mm}
\addtobeamertemplate{frametitle}{}{\vspace{-0.5em}}

\setbeamertemplate{navigation symbols}{}
\setbeamertemplate{footline}{%
  \hspace{5mm}{\usebeamerfont{page number in head/foot}\color{black!45}%
  ECON 282E \,$\cdot$\, \insertshortdeck}\hfill%
  \usebeamercolor[fg]{page number in head/foot}%
  \usebeamerfont{page number in head/foot}%
  \insertframenumber\,/\,\inserttotalframenumber\kern1em\vskip2pt%
}

% ---------------------------------------------------------------------------
% Chinese text under pdflatex (book titles etc.)
\newcommand{\zh}[1]{\begin{CJK*}{UTF8}{gbsn}#1\end{CJK*}}

% ---------------------------------------------------------------------------
% Deck title block
%   \decktitle{session}{deck id}{title}{date}
\newcommand{\insertshortdeck}{}
\newcommand{\decksession}{}
\newcommand{\deckid}{}
\newcommand{\decktitletext}{}
\newcommand{\deckdate}{}
\newcommand{\decktitle}[4]{%
  \renewcommand{\decksession}{#1}%
  \renewcommand{\deckid}{#2}%
  \renewcommand{\decktitletext}{#3}%
  \renewcommand{\deckdate}{#4}%
  \renewcommand{\insertshortdeck}{Session #1 \,$\cdot$\, Deck #1.#2}%
  \title{#3}%
  \author{Zhigang Feng}%
  \date{#4}%
}
\newcommand{\makedecktitle}{%
  \begin{frame}[plain,noframenumbering]
    \begin{tikzpicture}[remember picture,overlay]
      \fill[UCRBlue] (current page.north west) rectangle ([yshift=-0.55cm]current page.north east);
      \fill[UCRGold] ([yshift=-0.55cm]current page.north west) rectangle ([yshift=-0.62cm]current page.north east);
      \node[anchor=west, text=white, font=\small\bfseries] at ([xshift=5mm,yshift=-0.275cm]current page.north west)
        {ECON 282E \,$\cdot$\, Foundations of Macroeconomics};
      \node[anchor=east, text=white, font=\small] at ([xshift=-5mm,yshift=-0.275cm]current page.north east)
        {UC Riverside \,$\cdot$\, Fall 2026};
    \end{tikzpicture}
    \vspace*{1.1cm}
    {\usebeamerfont{title}\color{black!55}\large Session \decksession\ \,$\cdot$\, Deck \decksession.\deckid\par}
    \vspace{0.25cm}
    {\usebeamerfont{title}\color{RedTitle}\LARGE\bfseries \decktitletext\par}
    \vspace{0.7cm}
    {\large Zhigang Feng\par}
    \vspace{0.15cm}
    {\color{black!65}\deckdate\par}
    \vfill
    {\footnotesize\itshape\color{black!55}Build it on paper $\;\to\;$ solve it on the machine $\;\to\;$ push it to the frontier with AI\par}
  \end{frame}}

% ---------------------------------------------------------------------------
% Section divider (plain frame, not counted)
\newcommand{\sectiondivider}[2]{%
\begin{frame}[plain,noframenumbering]
\begin{center}
\vfill
{\Large\textbf{#1}\par}
\vspace{0.35cm}
{\normalsize #2\par}
\vfill
\end{center}
\end{frame}
}

% ---------------------------------------------------------------------------
% Boxes
\newtcolorbox{conceptbox}[1][]{colback=blue!3, colframe=RedTitle!55, boxrule=0.35mm,
  arc=1mm, left=2mm, right=2mm, top=1mm, bottom=1mm, #1}
\newtcolorbox{cautionbox}[1][]{colback=red!3, colframe=red!50!black, boxrule=0.35mm,
  arc=1mm, left=2mm, right=2mm, top=1mm, bottom=1mm, #1}
\newtcolorbox{examplebox}[1][]{colback=green!3, colframe=green!45!black, boxrule=0.35mm,
  arc=1mm, left=2mm, right=2mm, top=1mm, bottom=1mm, #1}
\newtcolorbox{takeawaybox}[1][]{colback=VeryLightGray, colframe=RedTitle!55, boxrule=0.35mm,
  arc=1mm, left=2mm, right=2mm, top=1mm, bottom=1mm, #1}
\newtcolorbox{researchbox}[1][]{enhanced, colback=violet!4, colframe=violet!60!black,
  boxrule=0.35mm, arc=1mm, left=2mm, right=2mm, top=1mm, bottom=1mm,
  fonttitle=\bfseries\small, title={Research in the wild}, #1}

\newcommand{\compacttables}{\renewcommand{\arraystretch}{1.15}}
\newcommand{\src}[1]{{\tiny\color{black!55}#1}}

% ---------------------------------------------------------------------------
%% Course map: the recurring "you are here" slide for ECON 282E.
%% Loaded by course_preamble.tex. Pattern follows AI-ECON-2026 lec01_map.tex, but the map
%% shows the whole ten-session course rather than one lecture.
%%
%%   \CourseMap{s}{label}   s = 1..10 highlights session s and prints "you are here: label"
%%                          (e.g. \CourseMap{1}{1.A2}); earlier sessions get a check mark.
%%   \CourseMap{0}{}        full overview, nothing highlighted.
%%
%% Note: TikZ option lists are comma-split before TeX conditionals run, so every \ifnum
%% below sits outside [...] option lists.

\newcommand{\CourseMap}[2]{%
\begin{center}
\begin{tikzpicture}[
    sess/.style={rectangle, rounded corners=2pt, draw=black!45, minimum width=1.34cm,
        minimum height=1.12cm, text width=1.26cm, align=center, inner sep=1pt},
    stage/.style={font=\tiny\bfseries, text=black!70},
]
  \foreach \i/\d/\t/\c in {%
      1/Sep 24/Intro \\ Growth/StageTrunk,
      2/Oct 1/HA $\cdot$ OLG \\ Policy/StageBranch,
      3/Oct 8/Num.\ DP \\ Python/StageMachine,
      4/Oct 15/Numerics \\ I/StageMachine,
      5/Oct 22/Numerics \\ II/StageMachine,
      6/Oct 29/The wall \\ \textbf{Midterm}/StageWall,
      7/Nov 5/ML \\ Deep nets/StageTools,
      8/Nov 12/RL \\ HA $+$ DL/StageTools,
      9/Nov 19/Text \\ LLMs/StageData,
      10/Dec 3/Agents \\ \textbf{Projects}/StageAgents} {
    \node[sess, fill=\c] (n\i) at ({1.46*(\i-1)}, 0)
      {{\scriptsize\bfseries S\i}\\[-1pt]{\tiny\color{black!60}\d}\\[1pt]{\tiny \t}};
    \ifnum\i<#1
      \node[font=\scriptsize, text=green!45!black] at ([xshift=-2.5mm,yshift=-2.2mm]n\i.north east) {$\checkmark$};
    \fi
    \ifnum\i=#1
      \draw[RedTitle, very thick, rounded corners=2pt]
        ([xshift=-0.6mm,yshift=-0.6mm]n\i.south west) rectangle ([xshift=0.6mm,yshift=0.6mm]n\i.north east);
      % keep the label inside the picture at both ends of the row
      \ifnum\i<3
        \node[font=\scriptsize\bfseries, text=RedTitle, anchor=south west] at ([yshift=1.2mm]n\i.north west)
          {$\blacktriangledown$\,you are here: #2};
      \else\ifnum\i>8
        \node[font=\scriptsize\bfseries, text=RedTitle, anchor=south east] at ([yshift=1.2mm]n\i.north east)
          {you are here: #2\,$\blacktriangledown$};
      \else
        \node[font=\scriptsize\bfseries, text=RedTitle, anchor=south] at ([yshift=1.2mm]n\i.north)
          {you are here: #2\,$\blacktriangledown$};
      \fi\fi
    \fi
  }
  % stage band under the sessions
  \foreach \a/\b/\lab/\c in {1/1/trunk/StageTrunk, 2/2/branches/StageBranch,
      3/5/the machine $\cdot$ classical methods/StageMachine, 6/6/the wall/StageWall,
      7/8/new tools/StageTools, 9/9/new data/StageData, 10/10/colleagues/StageAgents} {
    \fill[\c] ([yshift=-1.5mm]n\a.south west) rectangle ([yshift=-4.6mm]n\b.south east);
    \path ([yshift=-1.5mm]n\a.south west) -- ([yshift=-4.6mm]n\b.south east)
      node[midway, stage] {\lab};
  }
  \node[font=\footnotesize\itshape, text=black!70] at ([yshift=-9.5mm]$(n1.south)!0.5!(n10.south)$)
    {Build it on paper $\;\to\;$ solve it on the machine $\;\to\;$ push it to the frontier with AI};
\end{tikzpicture}
\end{center}
}


\graphicspath{{./figures/}}

\lstset{language=Python, basicstyle=\ttfamily\scriptsize, frame=none,
        backgroundcolor=\color{gray!8}, aboveskip=2pt, belowskip=2pt}

\decktitle{5}{B6}{Worked Example: The RBC Model by Projection}{Thursday, October 22, 2026}

\begin{document}

\makedecktitle

%=============================================================================
\begin{frame}{Where We Are}
\CourseMap{5}{5.B6}
\vspace{-1mm}
\begin{center}
\small \textbf{Block B:} 5.B1 the residual $\;\cdot\;$ 5.B2 choosing a basis $\;\cdot\;$ 5.B3 Chebyshev $\;\cdot\;$ 5.B4 finite elements $\;\cdot\;$ 5.B5 determining the coefficients $\;\cdot\;$ \textbf{5.B6} the worked example.
\end{center}
\end{frame}

%=============================================================================
\begin{frame}{The Same Model, a Fourth Way}
\begin{columns}[T]
\begin{column}{0.53\textwidth}
\small
Four sessions, one model.
\medskip
\begin{center}\footnotesize
\begin{tabularx}{\linewidth}{@{}>{\raggedright\arraybackslash}p{1.35cm} >{\raggedright\arraybackslash}X@{}}
\toprule
\textbf{3.A4} & grid VFI, everything crude \\
\textbf{4.B5} & grid VFI, Rouwenhorst $+$ interpolation $+$ golden section $+$ Howard \\
\textbf{5.A6} & perturbation around $k^{*}$ \\
\textbf{5.B6} & \textbf{Chebyshev collocation on a domain} \\
\bottomrule
\end{tabularx}
\end{center}
\medskip
Here the unknown is the \textbf{consumption policy}, written as a Chebyshev polynomial in $k$ for each of the seven shock states, and the coefficients are chosen to zero the Euler residual at the Chebyshev nodes.
\end{column}
\begin{column}{0.44\textwidth}
\begin{conceptbox}[title=\textbf{The commitments}]\small
\textbf{Basis:} Chebyshev in $k$, degree $7$ (5.B3).\\[1mm]
\textbf{Domain:} $k\in[0.5k^{*},1.5k^{*}]$ --- chosen, and it matters.\\[1mm]
\textbf{Shock:} the same Rouwenhorst $7$-state chain as 4.B5, so the comparison is like for like.\\[1mm]
\textbf{Projection:} orthogonal collocation (5.B5).
\end{conceptbox}
\end{column}
\end{columns}
\end{frame}

%=============================================================================
\begin{frame}[fragile]{Steps 1 and 2: The Basis, and the Residual}
\begin{columns}[T]
\begin{column}{0.57\textwidth}
\begin{lstlisting}
def psi(k):                       # [klo,khi] -> [-1,1]
    return 2*(k - klo)/(khi - klo) - 1

def cpol(th, k, j):               # policy for state j
    return chebval(psi(k), th[j])

def resid(th, k, j):
    c  = cpol(th, k, j)
    y  = z[j]*k**alpha + (1-delta)*k
    kp = np.clip(y - c, klo, khi)
    rhs = 0.0
    for l in range(n_z):
        cp = cpol(th, kp, l)
        R  = alpha*z[l]*kp**(alpha-1) + 1-delta
        rhs += Pi[j,l] * R/cp
    return 1.0 - c*beta*rhs
\end{lstlisting}
\end{column}
\begin{column}{0.40\textwidth}
\begin{conceptbox}[title=\textbf{Read the last line}]\small
That is the \textbf{unit-free Euler residual} of 3.A1 --- the same expression we have used to grade every solver since Session 3.
\medskip
The difference is that here it is not a diagnostic. It \emph{is} the equation being solved.
\end{conceptbox}
\medskip
\begin{cautionbox}\footnotesize
The \texttt{clip} matters: during the solve the iterate can propose a $k'$ outside the domain, where a Chebyshev polynomial diverges like $\cosh$.
\end{cautionbox}
\end{column}
\end{columns}
\end{frame}

%=============================================================================
\begin{frame}[fragile]{Step 3: Warm Start From Perturbation}
\begin{columns}[T]
\begin{column}{0.55\textwidth}
\begin{lstlisting}
# 5.A4's four numbers, reused as a starting guess
P, F, _, _ = klein(A, B, n_x=2)

nodes = cos(pi*(2*arange(1,n+2)-1)/(2*(n+1)))
kn = klo + 0.5*(nodes+1)*(khi - klo)

th0 = np.zeros((n_z, n+1))
for j in range(n_z):
    ch = F[0,0]*log(kn/kss) + F[0,1]*logz[j]
    th0[j] = chebfit(psi(kn), css*exp(ch), n)
\end{lstlisting}
\medskip
\small
The first-order perturbation policy, evaluated at the Chebyshev nodes and fitted. Six lines, and the nonlinear solve that follows converges from it without trouble.
\end{column}
\begin{column}{0.42\textwidth}
\begin{takeawaybox}\footnotesize
\textbf{This is the honest reason to learn Block A.} Projection needs a good initial guess for a $56$-dimensional nonlinear system, and perturbation produces one in under a millisecond.
\end{takeawaybox}
\medskip
\begin{examplebox}\footnotesize
5.B1 promised the two methods mix rather than compete. This is where it happens, and it is standard practice, not a trick.
\end{examplebox}
\end{column}
\end{columns}
\vspace{1mm}
\src{Ledger \texttt{worked\_proj.warm\_start}.}
\end{frame}

%=============================================================================
\begin{frame}[fragile]{Steps 4 and 5: Solve, Then Grade}
\begin{columns}[T]
\begin{column}{0.55\textwidth}
\begin{lstlisting}
def equations(flat):
    th = flat.reshape(n_z, n+1)
    return np.concatenate([resid(th, k_cheb, j)
                           for j in range(n_z)])

sol = root(equations, th0.ravel(),
           method="hybr", tol=1e-12)
th = sol.x.reshape(n_z, n+1)
#  56 unknowns, 0.13 s, residual 5.1e-15

E = euler_errors(cpol, th, n_test=2000)
#  graded on 2000 RANDOM points, not the nodes
\end{lstlisting}
\end{column}
\begin{column}{0.42\textwidth}
\begin{conceptbox}[title=\textbf{The grading rule}]\small
The residual is zero \emph{at the collocation nodes} by construction, so evaluating there would report $10^{-15}$ and mean nothing.
\medskip
Grade on \textbf{different points} --- 3.A1's separation of grids, and the single most common way to fool yourself here.
\end{conceptbox}
\medskip
\begin{examplebox}\footnotesize
$56$ unknowns: $8$ coefficients $\times$ $7$ shock states.
\end{examplebox}
\end{column}
\end{columns}
\vspace{1mm}
\src{Ledger \texttt{worked\_proj}: degree $7$, $56$ unknowns, $0.129$\,s.}
\end{frame}

%=============================================================================
\begin{frame}{Does It Work? Plot the Saving, Not the Policy}
\begin{columns}[T]
\begin{column}{0.53\textwidth}
\begin{center}
\begin{tikzpicture}
\begin{axis}[
    width=7.7cm, height=5.3cm,
    xmin=14, xmax=43, ymin=-0.62, ymax=0.62,
    xlabel={\scriptsize capital $k$}, ylabel={\scriptsize saving $g(k)-k$},
    tick label style={font=\scriptsize}, label style={font=\scriptsize},
    scaled ticks=false, /pgf/number format/fixed,
    /pgf/number format/precision=1,
    axis x line*=bottom, axis y line*=left]
  \addplot[very thick, RedTitle] coordinates {
      (14.1742,0.49434) (14.7431,0.47710) (15.3119,0.45952) (15.8808,0.44164) (16.4497,0.42346) 
      (17.0185,0.40502) (17.5874,0.38632) (18.1563,0.36739) (18.7251,0.34823) (19.2940,0.32887) 
      (19.8629,0.30931) (20.4317,0.28957) (21.0006,0.26964) (21.5694,0.24956) (22.1383,0.22931) 
      (22.7072,0.20892) (23.2760,0.18838) (23.8449,0.16771) (24.4138,0.14691) (24.9826,0.12599) 
      (25.5515,0.10496) (26.1204,0.08381) (26.6892,0.06256) (27.2581,0.04120) (27.8270,0.01975) 
      (28.3958,-0.00180) (28.9647,-0.02343) (29.5336,-0.04515) (30.1024,-0.06695) 
      (30.6713,-0.08883) (31.2401,-0.11078) (31.8090,-0.13281) (32.3779,-0.15491) 
      (32.9467,-0.17707) (33.5156,-0.19930) (34.0845,-0.22160) (34.6533,-0.24395) 
      (35.2222,-0.26636) (35.7911,-0.28883) (36.3599,-0.31136) (36.9288,-0.33393) 
      (37.4977,-0.35656) (38.0665,-0.37924) (38.6354,-0.40196) (39.2043,-0.42473) 
      (39.7731,-0.44754) (40.3420,-0.47040) (40.9108,-0.49330) (41.4797,-0.51624) 
      (42.0486,-0.53921) };
  \draw[black!45] (axis cs:14,0) -- (axis cs:43,0);
  \draw[black!35, dotted, thick] (axis cs:28.348,-0.62) -- (axis cs:28.348,0.62);
  \node[font=\tiny, text=black!55, anchor=south west] at (axis cs:28.6,-0.58) {$k^{*}=28.35$};
\end{axis}
\end{tikzpicture}
\end{center}
\end{column}
\begin{column}{0.44\textwidth}
\small
Saving $g(k)-k$ at the middle shock state. Plotting $g(k)$ itself would show a straight line at $45^{\circ}$ and hide everything.
\medskip

The curve crosses zero \textbf{at $k^{*}$} --- saving there is $2.8\times10^{-6}$ --- is positive below and negative above, and is smooth throughout. Three behaviour checks on one picture.
\medskip

Compare 4.B5, where the same crossing landed \textbf{ten grid spacings} away from $k^{*}$ because the policy was quantized. Here there is no grid to quantize it.
\end{column}
\end{columns}
\vspace{1mm}
\src{Ledger \texttt{fig\_proj\_policy}; the deviation is formed by pgfplots from the ledger's $(k,g)$ pairs.}
\end{frame}

%=============================================================================
\begin{frame}{What One More Coefficient Buys}
\begin{columns}[T]
\begin{column}{0.5\textwidth}
\small
Re-solve at a range of degrees, everything else identical:
\medskip
\begin{center}\footnotesize
\begin{tabularx}{\linewidth}{@{}r r >{\raggedleft\arraybackslash}X r@{}}
\toprule
degree & unknowns & $\max|\mathcal{E}|$ & sec \\
\midrule
$3$ & $28$ & $10^{-3.84}$ & $0.06$ \\
$5$ & $42$ & $10^{-5.15}$ & $0.09$ \\
$7$ & $56$ & $10^{-6.40}$ & $0.13$ \\
$9$ & $70$ & $10^{-7.63}$ & $0.18$ \\
$11$ & $84$ & $\mathbf{10^{-8.85}}$ & $0.50$ \\
\bottomrule
\end{tabularx}
\end{center}
\medskip
\textbf{About $1.2$ decades for every two degrees}, dead straight over five decades of error. That is the geometric convergence 5.B3 promised for a smooth function, arriving on schedule.
\end{column}
\begin{column}{0.47\textwidth}
\begin{takeawaybox}\footnotesize
Compare a grid method, where 3.A1 measured the value error falling like $h^{2}$ and the policy error like $h$. To gain $1.2$ decades there you multiply the grid by $16$. Here you add two coefficients.
\end{takeawaybox}
\medskip
\begin{cautionbox}\footnotesize
This only holds while the policy is \textbf{smooth}. Put a binding constraint in the domain and the same table would flatten out at $O(1/n)$ --- 5.B2's kinked function, exactly.
\end{cautionbox}
\end{column}
\end{columns}
\vspace{1mm}
\src{Ledger \texttt{worked\_proj.rows}; all five converged, graded on the same $2{,}000$ random points.}
\end{frame}

%=============================================================================
\begin{frame}{One Model, Four Methods}
\small
\begin{center}
\begin{tabularx}{0.97\linewidth}{@{}>{\raggedright\arraybackslash}p{3.3cm} r r >{\raggedleft\arraybackslash}X >{\raggedleft\arraybackslash}X@{}}
\toprule
\textbf{Method} & \textbf{numbers} & \textbf{sec} & $\max|\mathcal{E}|$ \textbf{near} $k^{*}$ & $\max|\mathcal{E}|$ \textbf{wide} \\
\midrule
global VFI (4.B5) & $2{,}100$ & $0.94$ & $10^{-2.82}$ & $10^{-2.51}$ \\
perturbation, 1st order & $\mathbf{4}$ & $<0.001$ & $10^{-4.01}$ & $10^{-2.56}$ \\
perturbation, 2nd order & $12$ & $0.02$ & $10^{-6.14}$ & $10^{-4.50}$ \\
Chebyshev collocation & $56$ & $0.13$ & $10^{-6.82}$ & $\mathbf{10^{-6.40}}$ \\
\bottomrule
\end{tabularx}
\end{center}
\medskip
\begin{columns}[T]
\begin{column}{0.49\textwidth}
\begin{takeawaybox}\footnotesize
\textbf{The grid method is last on every measure.} It stores $37\times$ more numbers than collocation, takes $7\times$ longer, and is \textbf{four decades} less accurate.
\end{takeawaybox}
\end{column}
\begin{column}{0.47\textwidth}
\begin{conceptbox}[title=\textbf{The flat column}]\footnotesize
Only collocation is nearly as accurate far from $k^{*}$ as near it: $10^{-6.40}$ against $10^{-6.82}$. That is what ``global'' means, measured.
\end{conceptbox}
\end{column}
\end{columns}
\vspace{1mm}
\src{Ledger \texttt{three\_methods}. Identical model, chain, domain and residual definition; the 4.B5 row reproduces that deck's own ledger to $0.002$ decades and its $66$ policy updates exactly.}
\end{frame}

%=============================================================================
\begin{frame}{So Why Did We Spend Two Sessions on Grids?}
\begin{columns}[T]
\begin{column}{0.53\textwidth}
\small
Because the table is about \textbf{this model}: smooth, one-dimensional and unconstrained --- the best case for a polynomial basis and the worst for a grid. Change any one and the ranking moves.
\medskip
\begin{itemize}\small
  \item \textbf{A binding constraint.} The policy kinks; Chebyshev gets Gibbs, perturbation does not exist, and VFI does not care.
  \item \textbf{A discrete choice.} Same, worse.
  \item \textbf{More dimensions.} The grid dies first, then complete polynomials, then Smolyak.
\end{itemize}
\end{column}
\begin{column}{0.44\textwidth}
\begin{conceptbox}[title=\textbf{The decision}]\footnotesize
\textbf{Smooth, low-$d$, want speed} $\to$ perturbation.\\[0.5mm]
\textbf{Smooth, need global accuracy} $\to$ projection.\\[0.5mm]
\textbf{Kinks, constraints, discrete choice} $\to$ VFI or EGM.\\[0.5mm]
\textbf{High-$d$ and smooth} $\to$ Smolyak, then neural nets (S7).
\end{conceptbox}
\medskip
\begin{cautionbox}\footnotesize
And whichever you choose, grade it on points you did not solve at.
\end{cautionbox}
\end{column}
\end{columns}
\vspace{1mm}
\src{Comparisons of this kind are often attributed to Aruoba, Fern\'andez-Villaverde \& Rubio-Ram\'irez (2006), which compares VFI, perturbation, Chebyshev projection, finite elements and PEA on this model. The table above is our own run, on our own calibration.}
\end{frame}

%=============================================================================
\begin{frame}{What We Deliberately Did Not Do}
\begin{columns}[T]
\begin{column}{0.53\textwidth}
\small
\begin{itemize}\small
  \item \textbf{No constraint.} The one case where the grid methods win, and the case every model in Session 2 has.
  \item \textbf{No second endogenous state.} At $d=2$ the tensor basis is $8^{2}=64$ coefficients per shock state; at $d=5$ it is $32{,}768$. That is why 5.B3 spent three frames on Smolyak.
  \item \textbf{We chose the domain by hand.} Too narrow and the solution is wrong where you need it; too wide and the polynomial wastes degrees. Nothing in the method chooses it.
  \item \textbf{No guarantee off the domain.} Outside $[k_{lo},k_{hi}]$ the approximation does not degrade gracefully --- it diverges.
\end{itemize}
\end{column}
\begin{column}{0.44\textwidth}
\begin{cautionbox}[title=\textbf{The honest summary}]\small
Projection is the most accurate method in this course so far, on the problems it suits, and it has \textbf{two free choices} --- basis and domain --- that it cannot make for you and will not warn you about.
\end{cautionbox}
\medskip
\begin{examplebox}\footnotesize
Perturbation had one such choice (the expansion point) and it was forced. Projection has two and neither is.
\end{examplebox}
\end{column}
\end{columns}
\end{frame}

%=============================================================================
\begin{frame}{Takeaways}
\begin{takeawaybox}
\begin{itemize}\small
  \item The same model, a fourth way: an $8$-term Chebyshev polynomial per shock state, coefficients chosen to \textbf{zero the Euler residual at the Chebyshev nodes}. $56$ unknowns, $0.13$ seconds.
  \item \textbf{Warm-start from perturbation.} Block A's four numbers give the nonlinear solver its initial guess --- the two methods cooperate rather than compete.
  \item \textbf{Grade on points you did not solve at.} At the collocation nodes the residual is $10^{-15}$ by construction and means nothing.
  \item Measured convergence: \textbf{$1.2$ decades for every two degrees}, straight over five decades. That is the geometric rate smoothness buys.
  \item \textbf{One model, four methods}: the grid solver stores $37\times$ more numbers, runs $7\times$ slower and is \textbf{four decades} less accurate than collocation --- and only collocation is as accurate far from $k^{*}$ as near it.
  \item That ranking is a property of a \textbf{smooth, one-dimensional, unconstrained} model. Add a binding constraint and it reverses, which is the single most important sentence in this session.
\end{itemize}
\end{takeawaybox}
\end{frame}

%=============================================================================
\begin{frame}{Bridge: The Models We Actually Care About}
\begin{columns}[T]
\begin{column}{0.53\textwidth}
\small
Session 5 ends with the best solution of the growth model this course will produce by classical means. It also ends with the condition that made it possible, stated four times: \textbf{smooth, one-dimensional, unconstrained}.
\medskip

Session 2's models are none of those. Aiyagari households sit at a borrowing limit; Krusell--Smith carries a distribution as a state; the life-cycle model has eighty cohorts. Every method in Sessions 3 to 5 has now met the case that defeats it.
\medskip

Next week we take these tools to those models, and find out exactly where they stop.
\end{column}
\begin{column}{0.44\textwidth}
\begin{conceptbox}[title=\textbf{Next week (S6)}]\small
Heterogeneous agents by classical methods: \textbf{Aiyagari} as a nested fixed point, \textbf{Krusell--Smith}, \textbf{OLG} and transitions --- and the complexity wall that Sessions 7 and 8 exist to break.
\end{conceptbox}
\medskip
\begin{examplebox}\footnotesize
\textbf{Homework 1} was due today. The \textbf{practice midterm} is posted this afternoon; the \textbf{midterm} is next Thursday, Block B, closed book, covering Sessions 1--5.
\end{examplebox}
\end{column}
\end{columns}
\end{frame}

\end{document}
